\section{Stage 4: What Does the Instrument Measure?}
\label{sec:measure}

\subsection{Labels across the sub-questions of one answer (E24)}
\label{sec:relation}
\FigureRelation

\step{Question.}
If evasion were a property of an answer's style, one answer would get one label whatever it was asked. Answers shared by several sub-questions test this (CPU only, existing outputs).

\step{Result.}
Of the 3,079 training pairs of siblings (two sub-questions, one answer), 61.5\% carry different evasion types and 40.6\% different clarity levels.
This is not annotator noise: on dev, moving to a sibling's sub-question changes the label almost as often as moving to an unrelated exchange, and far more often than two annotators disagree on the same item (Figure~\ref{fig:relation}a).
The label belongs to the pair, not to the answer.
M1 already behaves accordingly, repeating its label on only 19\% of the 58 dev sibling pairs with disjoint annotator sets.

\step{Relevance is not answerhood.}
Political scientists score answer quality by how specific an answer is to its question \citep{alvarez-morrier-2025-measuring}. We computed this with the frozen Qwen3-8B as the gain in the answer's log-likelihood under its own sub-question over three random ones.
It is ordered across the nine labels (Kruskal--Wallis $p = 2 \times 10^{-12}$) but barely discriminative (AUROC 0.566 for Clear Reply against the rest): Partial, Implicit, Deflection and General answers fit their own question as well as Explicit ones do (Figure~\ref{fig:relation}b).
An answer can be \emph{about} the question without \emph{answering} it, which is the distinction the taxonomy draws.
Five further module ideas (epistemic-uncertainty routing, joint decoding of siblings, a parameter-free metric-optimal decision rule, single-order debiasing, conformal candidate sets) were tested on existing outputs and rejected before any new run (Appendix~\ref{sec:ledger}).

\subsection{Does the instrument read the question causally? (E25a)}

\step{Registered prediction (H25a-1).}
An answer paired with a question from another interview has, by the taxonomy's definition, not replied to it.
If the instrument models the relation, the Clear-Reply share should fall by at least half and at most 50\% of items should keep their label. If it labels answer style, labels should persist.
This is not obvious: training never showed a mismatched pair.

\step{Result.}
Held (Figure~\ref{fig:swap} in the appendix, left).
With an unrelated question the Clear-Reply share falls from 0.253 to 0.013, only 18.8\% of items keep their label, and 62 of the 78 Clear Replies become Dodging. The label-free replication on test runs in the same direction (0.338 to 0.110).
Without any question the share \emph{rises} to 0.640: with nothing asked, nothing can be evaded.
The instrument models the relation, the precondition for using it on new speakers.

\subsection{Does the question or the speaker drive evasion? (E25b)}

\step{Design.}
If evasion is relational, it could be driven by the question (some questions are hard to answer for anyone) or by the speaker (Presidents evade strategically).
We put the 308 dev sub-questions (191 interview turns) to two open instruction-tuned LLMs, Mistral-7B-Instruct-v0.3 \citep{jiang2023mistral} and Qwen3-8B with thinking off \citep{qwen3}, in four conditions that each differ from P in one factor. \textbf{N} is a neutral analyst with no stake, \textbf{P} the President in that interview on its date, \textbf{F} P followed by ``With respect, you haven't answered the question'', and \textbf{B} P with the President's real answer as private notes.
The frozen instrument labelled all 2,464 answers against their sub-questions. Statistics use a cluster bootstrap over turns and Mantel--Haenszel odds ratios pooled over models.

\FigureAudit

\begin{table}[t]
\centering
\small
\setlength{\tabcolsep}{3pt}
\begin{tabular}{@{}llll@{}}
\toprule
 & Registered prediction & Result & Verdict \\
\midrule
a-1 & unrelated question: CR $\le \tfrac12$ & 0.253 $\rightarrow$ 0.013 & held \\
b-1 & no stake: OR $\ge 1.5$ & 1.18 & failed \\
b-2 & stake: P $-$ N $\ge +0.10$ & $-$0.031 & failed \\
b-3 & pursuit: $\ge$50\% replies & 66\% & held \\
b-4 & notes: indirect $-$ non $\ge +0.15$ & +0.113 & failed \\
c & half answer: $\ge$70\% agree & 66\% & failed \\
d-1 & frozen probe: \Stwo{} $\ge 0.45$ & 0.397 & failed \\
d-2 & probe: siblings same $\ge$30\% & 45\% & held \\
\bottomrule
\end{tabular}
\caption{E25's registered hypotheses (dev, frozen instrument). CR is the Clear-Reply share and OR the odds ratio, shown beside each pair in Figure~\ref{fig:audit}b.}
\label{tab:e25}
\end{table}

\step{Results.}
Both registered hypotheses about the question and the stake failed (Table~\ref{tab:e25}, Figure~\ref{fig:audit}): a neutral analyst does not evade more on the questions the Presidents evaded (OR 1.18), and playing the President does not raise the evasion \emph{rate} ($-$0.031).
In the President's role, though, the same models evade on the questions the Presidents evaded (OR 2.74 \ci{1.87}{4.17}, with 3.0 for Mistral and 2.5 for Qwen). We did not register this comparison, so it is exploratory.
In plain rates (Figure~\ref{fig:audit}b), the neutral analyst evades on 56\% of the questions the Presidents evaded and on 52\% of those they answered. In the President's role the models evade on 58\% and 34\%.
The move therefore comes from the questions the Presidents answered directly, which the models now also answer directly, not from the questions the Presidents dodged. A neutral analyst addressed as ``Mr.\ President'' evades about half of everything, because it is not the addressee.
The role changes \emph{where} models evade, not how often (52\% against 55\%). A question's answerability shows only for a speaker who can answer it, and our data cannot tell this apart from strategic evasion.
Labelling N and P with one instrument removes a question prior that acts alike on every answer, not one that interacts with answer style, and the question-only control was degenerate.
Pursuit works: a follow-up turns 66\% of evasive answers into clear replies.
With the President's answer as notes, indirect answers become replies more often than non-answers, by less than the registered margin. Condition effects survive a length covariate, but the role odds ratio was not length-adjusted.

\step{Is the instrument valid on LLM text?}
The registered human check was dropped for lack of time. Without labels, the instrument is \emph{more} confident on LLM answers than on the Presidents' (seed agreement 0.81 vs 0.73), though confident errors remain possible.

\subsection{Is there a shortcut? (E25c--d)}
A shortcut would exist if evasion were visible early in an answer or encoded in the frozen base model. Neither holds (Figure~\ref{fig:swap}b--c).
Half an answer agrees with the full-answer label on only 66\% of items: refusals show early, General and Dodging only at the end.
Linear probes on the frozen Qwen3-8B-Base reach 0.397 \Stwo{} at best, against 0.525 for one fine-tuned seed, and they give disjoint siblings the same label 45\% of the time, against 19\% for M1. Fine-tuning is what ties the label to the question.

\step{What it raised.}
Measuring evasion needs the whole answer and the fine-tuned ten-seed ensemble, so one lever that leaves the instrument's model, and so its labels, unchanged is numeric precision.
